% ============================================================================
%  Subdocumento 6. Contenido. Se usa igual en el documento único y en el
%  PDF propio del subdocumento. Los formularios que le asigna el T-21 se
%  agregan solos al final (datos en formularios/ de esta carpeta).
% ============================================================================

\begin{resumenApertura}
\marcador{Qué resuelve este subdocumento, en cuatro o cinco líneas}
\begin{recibe}
  \item \marcador{Qué recibe el mandante}
\end{recibe}
\end{resumenApertura}

\section{Metodología de gestión del proyecto}\label{sec:6-metodologia-de-gestion-del-proye}

\marcador{PMBOK adaptado e integración de enfoques ágiles. El Formulario T-9 indica dónde está cada exigencia}

\section{Metodología de desarrollo}\label{sec:6-metodologia-de-desarrollo}

\marcador{Enfoque de desarrollo y sus implicancias. El Formulario T-10 indica dónde está cada exigencia}

\section{DevSecOps, integración y entrega continuas}\label{sec:6-devsecops-integracion-y-entrega-}

\marcador{Infraestructura como código y automatización de pruebas}

\section{Ceremonias, artefactos, cadencias y decisiones}\label{sec:6-ceremonias-artefactos-cadencias-}

\marcador{Mecanismos de decisión del proyecto}
